% Part: turing-machines
% Chapter: undecidability
% Section: representing-tms

\documentclass[../../../include/open-logic-section]{subfiles}

\begin{document}

\olfileid{tur}{und}{rep}
\olsection{ट्यूरिंग यंत्रांचे निरूपण}

\begin{explain}
ट्यूरिंग यंत्रे आणि त्यांचे वर्तन प्रथम-क्रम तर्कशास्त्रातील
!!a{sentence} वाक्याने निरूपित करण्यासाठी योग्य भाषा
परिभाषित करावी लागते. त्या भाषेचे दोन भाग आहेत:
यंत्राच्या स्थितिवर्णनांचे वर्णन करणारी !!{predicate}s
म्हणजे विधेये, आणि संगणनाच्या पायऱ्या (“क्षण”)
तसेच फितीवरील स्थाने यांना संख्या देण्यासाठीचे
पदबंध.

दोन प्रकारची !!{predicate}s म्हणजे विधेये
आपण घेऊ; दोन्ही द्विस्थानी आहेत: प्रत्येक
अवस्था~$q$ साठी !!a{predicate}~$\Obj Q_q$,
आणि फितीवरील प्रत्येक चिन्ह~$\sigma$
साठी !!a{predicate}~$\Obj S_\sigma$.
पहिल्या प्रकाराने~$M$ ची अवस्था आणि
त्याच्या फितीवरील शीर्षाचे स्थान वर्णन
करता येते; दुसऱ्या प्रकाराने फितीवरील
मजकूर वर्णन करता येतो.

फितीवरील शीर्षाचे स्थान आणि केलेल्या
पायऱ्यांची संख्या व्यक्त करण्यासाठी
संख्या व्यक्त करण्याचा मार्ग हवा.
यासाठी !!a{constant}~$\Obj 0$
आणि उत्तरवर्ती फलन~$\prime$ हे
$1$-स्थानी फलन वापरतो. संकेताप्रमाणे
ते त्याच्या आदानाच्या \emph{नंतर}
लिहितो (आणि कंस गाळतो).

प्रत्येक संख्या $n$ साठी तिचे
~$\Lang L_M$ मधील निरूपण करणारे
मानक पद~$\num{n}$ असते; त्याला
~$n$ चा \emph{संख्यांक} म्हणतात.
$\num{0}$ म्हणजे $\Obj 0$,
$\num{1}$ म्हणजे $\Obj 0'$,
$\num{2}$ म्हणजे $\Obj 0''$,
इत्यादी. अधिक औपचारिकरीत्या:
\begin{align*}
\num{0} & = \Obj 0 \\
\num{n+1} &= \num{n}'
\end{align*}
$\num{0}$ हे पद, म्हणजे $\Obj 0$,
फितीवरील सर्वांत डावीकडील घराचे स्थान
आणि संगणनाची पहिली पायरी सुरू होण्यापूर्वीचा
क्षण (प्रारंभिक स्थितिवर्णन) या दोन्हींना
नाव देते. $\num{1}$ हे पद, म्हणजे
$\Obj 0'$, सर्वांत डावीकडील घराच्या
उजवीकडील घराला आणि संगणनाच्या
पहिल्या पायरीनंतरच्या क्षणाला नाव
देते; पुढेही याचप्रमाणे.

फितीवरील स्थानांचा क्रम (“याच्या डावीकडे”)
आणि संगणनाच्या पायऱ्यांचा क्रम (“याच्या
आधी”) हे दोन्ही व्यक्त करण्यासाठी
!!a{predicate}~$<$ सुद्धा घेऊ.

भाषा ठरल्यावर $!T(M,
w)$ ची “स्वयंसिद्धके” यादीत देऊ.
म्हणजे आदान~$w$ वर चालवलेल्या~$M$
चे वर्तन एकत्रितपणे वर्णन करणारी
!!{sentence}s वाक्ये. त्यांत
$\Obj 0$, $\prime$ आणि $<$ यांवर
अटी घालणारी !!{sentence}s,
आरंभीचे स्थितिवर्णन सांगणारी
!!{sentence}s, तसेच विशिष्ट सूचना
राबवल्यानंतर~$M$ चे स्थितिवर्णन
काय होते ते सांगणारी !!{sentence}s
असतील.
\end{explain}

\begin{defn}
  \ollabel{defn:tm-descr}
$M = \tuple{Q, \Sigma, q_0, \delta}$
हे ट्यूरिंग यंत्र दिले असता
भाषा~$\Lang L_M$ मध्ये पुढील
गोष्टी असतात:
\begin{enumerate}
\item प्रत्येक अवस्था~$q \in Q$ साठी
  द्विस्थानी !!{predicate} $\Obj Q_q(x, y)$.
  अंतर्ज्ञानाच्या दृष्टीने,
  $\Obj Q_q(\num{m}, \num{n})$
  याचा अर्थ “$n$ पायऱ्यांनंतर
  $M$ हे~$q$ अवस्थेत असून
  $m$-व्या घराकडे पाहत आहे.”
\item प्रत्येक चिन्ह~$\sigma\in \Sigma$
  साठी द्विस्थानी !!{predicate}
  $\Obj S_\sigma(x, y)$.
  अंतर्ज्ञानाच्या दृष्टीने,
  $\Obj S_\sigma(\num{m},
  \num{n})$ याचा अर्थ “$n$
  पायऱ्यांनंतर $m$-व्या घरात
  चिन्ह~$\sigma$ आहे.”
\item एक !!{constant} $\Obj 0$
\item एक एकस्थानी !!{function} $\prime$
\item एक द्विस्थानी !!{predicate} $<$
\end{enumerate}
\end{defn}

आदान $w = \sigma_{i_1}\dots\sigma_{i_k}$
वर ट्यूरिंग यंत्र~$M$ कसे चालते याचे
वर्णन करणारी !!{sentence}s वाक्ये
पुढीलप्रमाणे आहेत:
\begin{enumerate}
\item संख्या आणि~$<$ यांचे वर्णन
करणारी स्वयंसिद्धके:
\begin{enumerate}
\item प्रत्येक संख्या तिच्या उत्तरवर्तीपेक्षा
लहान आहे असे सांगणारे !!^a{sentence}:
\[
\lforall[x][x < x']
\]
\item $<$ संक्रमणशील आहे याची
खात्री करणारे !!^a{sentence}:
\[
\lforall[x][\lforall[y][\lforall[z][
      ((x < y \land y < z) \lif x < z)]]]
\]
\end{enumerate}
\item आदानाच्या प्रारंभिक स्थितिवर्णनाचे
वर्णन करणारी स्वयंसिद्धके:
\begin{enumerate}
\item $0$~पायऱ्यांनंतर---म्हणजे
  यंत्र सुरू होण्यापूर्वी---$M$
  हे आरंभीच्या अवस्था~$q_0$
  मध्ये असून घर~$1$ कडे पाहते:
\[
\Obj Q_{q_0}(\num{1}, \num{0})
\]
\item पहिल्या $k+1$ घरांत
  $\TMendtape$, $\sigma_{i_1}$,
  \dots, $\sigma_{i_k}$ ही चिन्हे
  असतात:
\[
\Obj S_\TMendtape(\num{0}, \num{0}) \land
\Obj S_{\sigma_{i_1}}(\num{1}, \num{0}) \land
\dots \land
\Obj S_{\sigma_{i_k}}(\num{k}, \num{0})
\]
\item बाकी फीत रिकामी असते:
\[
\lforall[x][(\num{k} < x \lif \Obj S_\TMblank(x, \num{0}))]
\]
\end{enumerate}
\item एका स्थितिवर्णनातून पुढच्या
स्थितिवर्णनाकडे जाणाऱ्या संक्रमणांचे
वर्णन करणारी स्वयंसिद्धके:

पुढे $!A(x, y)$ हे पुढील
रूपातील सर्व !!{sentence}s वाक्यांच्या
संयोगाने मिळते असे मानू:
\[
\lforall[z][
  (((z < x \lor x < z) \land \Obj S_\sigma(z, y))
  \lif \Obj S_\sigma(z, y'))]
\]
येथे दोन चल मुक्त असल्याने हा स्वतः स्वतंत्र वाक्य
नसून सूत्रांचा रूपबंध आहे; पुढील संक्रमण-वाक्यांत
हे चल सार्वत्रिक संख्यापकांखाली येतात.
%READERNOTE{OLTUR-010 — गोठवलेल्या इंग्रजी मजकुरात A(x,y) बनवणाऱ्या अटींना sentences म्हटले आहे. प्रदर्शित अटींमध्ये x आणि y मुक्त आहेत; म्हणून त्या स्वतंत्र वाक्ये नसून उघडी सूत्रे आहेत. पुढील संक्रमण-स्वयंसिद्धकांत त्या सार्वत्रिक संख्यापकांच्या व्याप्तीत वापरल्या जातात.}
येथे $\sigma \in \Sigma$.
$!A(\num{m},\num{n})$ वापरून
“घर~$m$ सोडल्यास, $n+1$
पायऱ्यांनंतरची फीत $n$
पायऱ्यांनंतर होती तशीच आहे”
हे व्यक्त करतो.
\begin{enumerate}
\item \ollabel{rep-right} प्रत्येक
  $\delta(q_i, \sigma) =
  \tuple{q_j, \sigma', \TMright}$
  सूचनेसाठी पुढील !!{sentence}:
\begin{align*}
& \lforall[x][\lforall[y][(
   (\Obj Q_{q_i}(x, y) \land \Obj S_{\sigma}(x, y)) \lif {}]] \\
&\qquad   (\Obj Q_{q_j}(x', y') \land \Obj S_{\sigma'}(x, y') \land
!A(x, y)))
\end{align*}
याचा अर्थ असा: $y$ पायऱ्यांनंतर
यंत्र~$q_i$ अवस्थेत असून
चिन्ह~$\sigma$ असलेल्या
घर~$x$ कडे पाहत असेल,
तर $y+1$ पायऱ्यांनंतर
ते घर~$x+1$ कडे पाहील,
~$q_j$ अवस्थेत असेल,
घर~$x$ मध्ये आता~$\sigma'$
असेल आणि घर~$x$ वगळता
इतर प्रत्येक घरात~$y$
पायऱ्यांनंतर होते तेच
चिन्ह असेल.

\item \ollabel{rep-left} प्रत्येक
  $\delta(q_i, \sigma) =
  \tuple{q_j, \sigma', \TMleft}$
  सूचनेसाठी पुढील !!{sentence}:
\begin{align*}
& \lforall[x][\lforall[y][
    ((\Obj Q_{q_i}(x', y) \land \Obj S_{\sigma}(x', y)) \lif {}]]\\
& \qquad   (\Obj Q_{q_j}(x, y') \land \Obj S_{\sigma'}(x', y') \land
!A(x', y))) \land {}\\
& \lforall[y][((\Obj Q_{q_i}(\num{0}, y) \land \Obj S_{\sigma}(\num{0},
    y)) \lif {}]\\
& \qquad (\Obj Q_{q_j}(\num{0}, y') \land \Obj S_{\sigma'}(\num{0},
  y') \land !A(\num{0}, y)))
\end{align*}
ही रचना नीट समजून घ्या:
आता आपण “घर~$x$ कडे
पाहत असेल तर~\dots”
याऐवजी “घर $x+1$
कडे पाहत असेल तर~\dots”
असे म्हणतो. डावीकडे
जाण्याचा अर्थ पुढच्या
पायरीत यंत्र घर~$x$
कडे पाहते. पण लिहिले
जाते ते घर~$x+1$.
वजाबाकी किंवा पूर्ववर्ती
फलन आपल्याकडे नसल्यामुळे
ही रचना वापरतो.
पहिल्या अटीतील स्थिरता-नियम
लिहिले जाणारे घरच
वगळतो; डावीकडे गेलेले
नवे शीर्षस्थान नाही.
%READERNOTE{OLTUR-009 — गोठवलेल्या इंग्रजी डावीकडील संक्रमणाच्या पहिल्या फलितात A(x,y) आहे. पण त्या संक्रमणात चिन्ह लिहिले जाते x' या घरात; A च्या आधीच्या व्याख्येनुसार त्याने लिहिले जाणारे घर वगळले पाहिजे. येथे A(x',y) केले आहे. दुसऱ्या, घर शून्यवरील शाखेतील A(0,y) बदललेले नाही.}

$x+1$ या रूपातील संख्या
$1$, $2$, \dots,
असतात हे लक्षात घ्या.
म्हणून घर~$0$ कडे पाहताना
डावीकडे जाण्याची सूचना
असल्याची बाब यात येत
नाही (आणि अर्थातच ते
डावीकडे जाऊ शकत नाही---त्याच
घरावर राहते). दुसऱ्या
संयुक्त अटीत ही विशेष
बाब येते: $y$ पायऱ्यांनंतर
यंत्र घर~$0$ कडे पाहत,
अवस्था $q_i$ मध्ये असेल
आणि घर~$0$ मध्ये
चिन्ह~$\sigma$ असेल, तर
$y+1$ पायऱ्यांनंतर ते
घर~$0$ कडेच पाहत असेल,
नव्या अवस्था~$q_j$ मध्ये
असेल, घर~$0$ वरील
चिन्ह $\sigma'$ असेल
आणि घर~$0$ वगळता
इतर घरांत $y$~पायऱ्यांनंतर
होती तीच चिन्हे असतील.
\item \ollabel{rep-stay} प्रत्येक
  $\delta(q_i, \sigma) =
  \tuple{q_j, \sigma', \TMstay}$
  सूचनेसाठी पुढील !!{sentence}:
\begin{align*}
& \lforall[x][\lforall[y][(
   (\Obj Q_{q_i}(x, y) \land \Obj S_{\sigma}(x, y)) \lif {}]] \\
&\qquad   (\Obj Q_{q_j}(x, y') \land \Obj S_{\sigma'}(x, y') \land
!A(x, y)))
\end{align*}
\end{enumerate}
\end{enumerate}
ट्यूरिंग यंत्र~$M$ आणि आदान~$w$
यांच्यासाठी वरील सर्व !!{sentence}s
वाक्यांच्या संयोगाला $!T(M, w)$ म्हणू.

~$M$ अखेरीस थांबते हे व्यक्त
करण्यासाठी “काही पायऱ्यांनंतर
संक्रमण फलन अपरिभाषित होईल”
असे सांगणारे !!{sentence}
वाक्य हवे. ज्या सर्व
$\tuple{q, \sigma}$ जोड्यांसाठी
~$\delta(q, \sigma)$ अपरिभाषित
आहे, त्यांचा संच~$X$
असू द्या. मग~$!E(M, w)$
हे पुढील !!{sentence} वाक्य
असू द्या:
\[
\lexists[x][\lexists[y][(\bigvee_{\tuple{q, \sigma} \in
      X}(\Obj Q_q(x, y) \land \Obj S_\sigma(x, y)))]]
\]

खास थांबण्याची अवस्था~$h$
असलेले ट्यूरिंग यंत्र
वापरले तर हे आणखी सोपे
होते: अशा वेळी पुढील
!!{sentence} वाक्य~$!E(M, w)$
यंत्र अखेरीस थांबते हे
व्यक्त करते:
\[
\lexists[x][\lexists[y][\Obj Q_h(x, y)]]
\]

\begin{prop}
\ollabel{prop:mlessk}
$m < k$ असेल, तर
$!T(M, w) \Entails \num{m} < \num{k}$.
\end{prop}

\begin{proof}
सराव.
\end{proof}

\begin{prob}
\olref[tur][und][rep]{prop:mlessk}
सिद्ध करा. (सूचना: $k-m$
वर गणितीय विगमन वापरा.)
\end{prob}

\end{document}
